% =====================================================
% 题型：直三棱柱外接球填空题 (q_solid_prism_sphere.tex)
% =====================================================
% =====================================================
% 难度：★★ (中档)
% =====================================================
\newcommand{\qSolidPrismSphereStem}{在直三棱柱 \(ABC-A_1B_1C_1\) 中，\(2AB = AA_1 = 2\)，若该三棱柱的外接球的表面积为 \(8\pi\)，则 \(\angle ACB\) 的大小为 \underline{\hspace{2.5cm}}。}

\newcommand{\qSolidPrismSphereFigure}[1][1.3]{%
  \begin{tikzpicture}[scale=#1, >=stealth]
    \coordinate (A) at (0,0);
    \coordinate (B) at (1.75,0);
    \coordinate (C) at (2.45,0.62);
    \coordinate (A1) at (0,1.85);
    \coordinate (B1) at (1.75,1.85);
    \coordinate (C1) at (2.45,2.47);

    \draw[dashed, thick, gray!75] (A) -- (C);

    \draw[thick] (A) -- (B) -- (C);
    \draw[thick] (A) -- (A1) -- (B1) -- (B);
    \draw[thick] (A1) -- (C1) -- (B1);
    \draw[thick] (C) -- (C1);

    \node[below left] at (A) {\small $A$};
    \node[below] at (B) {\small $B$};
    \node[right] at (C) {\small $C$};
    \node[left] at (A1) {\small $A_1$};
    \node[left] at (B1) {\small $B_1$};
    \node[right] at (C1) {\small $C_1$};
  \end{tikzpicture}%
}

\newcommand{\qSolidPrismSphereAnswer}{\(30^\circ\) 或 \(150^\circ\)}

\newcommand{\qSolidPrismSphereSolution}{%
  设直三棱柱底面外接圆的半径为 \(r\)，高为 \(h\)，外接球的半径为 \(R\)。
  因为直三棱柱的外接球与该三棱柱共用同一个球心与半径，
  根据外接球的几何性质，有：
  \(R^2 = r^2 + \left(\dfrac{h}{2}\right)^2\)。\\
  由已知，\(2AB = AA_1 = 2 \implies AB = 1\)，且三棱柱的高为 \(h = AA_1 = 2\)。
  因为外接球的表面积为 \(8\pi\)，
  所以 \(4\pi R^2 = 8\pi \implies R^2 = 2\)。
  将 \(R^2 = 2\)，\(h = 2\) 代入公式，可得：
  \(2 = r^2 + \left(\dfrac{2}{2}\right)^2 \implies 2 = r^2 + 1 \implies r^2 = 1 \implies r = 1\)。
  所以在底面 \(\triangle ABC\) 中，其外接圆半径为 \(r = 1\)。
  根据正弦定理，有：
  \(\dfrac{AB}{\sin\angle ACB} = 2r \implies \dfrac{1}{\sin\angle ACB} = 2 \implies \sin\angle ACB = \dfrac{1}{2}\)。
  在 \(\triangle ABC\) 中，\(\angle ACB \in (0, \pi)\)，
  所以 \(\angle ACB = 30^\circ\) 或 \(\angle ACB = 150^\circ\)。
  故填 \(30^\circ\) 或 \(150^\circ\)。%
}

\newcommand{\qSolidPrismSphereDiagnosis}{%
  若只写 \(30^\circ\)，则说明由 \(\sin C=\frac12\) 求三角形内角时漏掉了钝角解 \(150^\circ\)。%
}

\newcommand{\qSolidPrismSphereTeachingNotes}{%
  快讲。强调由正弦值求三角形内角时，必须检查 \(C\in(0,\pi)\) 下的双解。%
}
